
\documentclass{article}
\usepackage{colortbl}
\usepackage{makecell}
\usepackage{multirow}
\usepackage{supertabular}

\begin{document}

\newcounter{utterance}

\centering \large Interaction Transcript for game `clean\_up', experiment `2\_hard\_3obj\_hr', episode 2 with Qwen3.6{-}35B{-}A3B{-}FP8{-}without{-}reasoning.
\vspace{24pt}

{ \footnotesize  \setcounter{utterance}{1}
\setlength{\tabcolsep}{0pt}
\begin{supertabular}{c@{$\;$}|p{.15\linewidth}@{}p{.15\linewidth}p{.15\linewidth}p{.15\linewidth}p{.15\linewidth}p{.15\linewidth}}
    \# & $\;$A & \multicolumn{4}{c}{Game Master} & $\;\:$B\\
    \hline

    \theutterance \stepcounter{utterance}  
    & & \multicolumn{4}{p{0.6\linewidth}}{
        \cellcolor[rgb]{0.9,0.9,0.9}{
            \makecell[{{p{\linewidth}}}]{
                \texttt{\tiny{[P1$\langle$GM]}}
                \texttt{Ja sam vaš voditelj igre, a vi igrate suradničku igru sa sljedećom mrežom kao pločom za igru:} \\
\\ 
\texttt{```} \\
\texttt{ 1234567} \\
\texttt{╔═╤═╤═╤═╗} \\
\texttt{║◌│◌│◌│◌║ 1} \\
\texttt{╟─┼─┤◌│◌║ 2} \\
\texttt{║◌│◌│◌│◌║ 3} \\
\texttt{║L│◌├─┴─╢ 4} \\
\texttt{╟─┴─┤◌◌◌║ 5} \\
\texttt{║◌◌C└─┬─╢ 6} \\
\texttt{║◌◌◌◌P│◌║ 7} \\
\texttt{╚═════╧═╝} \\
\\ 
\texttt{```} \\
\\ 
\texttt{* Gornji rub prikazuje x{-}koordinate koje rastu udesno, a desni rub y{-}koordinate koje rastu prema dolje.} \\
\texttt{* Sljedeći objekti nasumično su postavljeni na vašoj mreži: 'C', 'L', 'P'.} \\
\\ 
\\ 
\texttt{Drugi igrač vidi varijantu ploče za igru, gdje su objekti postavljeni na različitim nasumičnim lokacijama. Ne možete vidjeti ploču drugog igrača, a ni on ne može vidjeti vašu.} \\
\\ 
\\ 
\texttt{**Cilj:**} \\
\\ 
\texttt{Oba igrača trebaju pomicati objekte na svojim pozadinama tako da identični objekti završe na istim koordinatama. Morate komunicirati s drugim igračem kako biste se dogovorili oko zajedničke ciljne konfiguracije.} \\
\\ 
\\ 
\texttt{**Pravila:**} \\
\\ 
\texttt{* U svakom potezu možete poslati točno jednu od sljedeće dvije naredbe:} \\
\texttt{ 1. `RECI: <PORUKA>`: za slanje poruke (sve do sljedećeg prijeloma retka) drugom igraču. Proslijedit ću je vašem partneru.} \\
\texttt{2. `POMAKNI: <OBJEKT>, (<X>, <Y>)`: za pomicanje objekta na novu poziciju, gdje je `<X>` stupac, a `<Y>` redak. Obavijestit ću vas je li vaš potez bio uspješan ili ne.} \\
\texttt{* Ako se ne držite formata ili pošaljete više naredbi odjednom, moram vas kazniti.} \\
\texttt{* Ako oba igrača skupe više od 8 kazni, oboje gubite igru.} \\
\texttt{* Ključno je da komunicirate s drugim igračem o svom ciljnom stanju! *Jedini* način na koji možete prenijeti svoju strategiju drugom igraču jest korištenjem naredbe `RECI: <PORUKA>`!} \\
\\ 
\\ 
\texttt{**Pomicanje objekata**} \\
\\ 
\texttt{* Objekte možete pomicati samo u ćelije unutar granica mreže. Ciljna ćelija mora biti prazna, tj. smije sadržavati samo simbol '◌'.} \\
\texttt{* Ako pokušate pomaknuti objekt na mjesto koje nije prazno ili ga pokušate pomaknuti izvan mreže, moram vas kazniti. Dobivate još jedan pokušaj.} \\
\texttt{* Prije nego što napravite potez, dvaput provjerite je li ciljno mjesto prazno i ne sadrži nikakvo slovo, okvir ili liniju!} \\
\\ 
\\ 
\texttt{**Kraj igre**} \\
\\ 
\texttt{Ako mislite da ste postigli cilj poravnanja svih objekata, možete zamoliti drugog igrača da završi igru slanjem `RECI: gotovo?`. Ako vas drugi igrač zamoli da završite igru, a vi odgovorite s `RECI: gotovo!`, igra će završiti.} \\
\\ 
\texttt{Oba igrača pobjeđuju ako igra završi unutar 12 rundi, pri čemu je jedna runda definirana kao to da oba igrača pošalju valjanu naredbu.} \\
\\ 
\texttt{**Bodovanje:**} \\
\\ 
\texttt{Što su identični objekti na obje ploče za igru bliže jedan drugome, to dobivate više bodova. Kazne smanjuju vaše bodove. Možete li srušiti rekord?} \\
\\ 
\\ 
\texttt{Molimo pošaljite poruku drugom igraču kako biste započeli igru!} \\
            }
        }
    }
    & \\ \\

    \theutterance \stepcounter{utterance}  
    & \multicolumn{4}{p{0.6\linewidth}}{
        \cellcolor[rgb]{0.9,0.9,0.9}{
            \makecell[{{p{\linewidth}}}]{
                \texttt{\tiny{[P1$\rangle$GM]}}
                \texttt{RECI: Pozdrav! Da bismo uskladili pozicije objekata C, L i P, dogovorimo se da ćemo ih sve premjestiti u gornji lijevi kut mreže. Cilj nam je da svi tri objekta završe na koordinati (1, 1). Molim te, provjeri gdje se nalazi tvoj objekt 'C' i premjesti ga na (1, 1). Ja ću učiniti isto sa svojim objektima.} \\
            }
        }
    }
    & & \\ \\

    \theutterance \stepcounter{utterance}  
    & & & \multicolumn{2}{p{0.3\linewidth}}{
        \cellcolor[rgb]{0.9,0.9,0.9}{
            \makecell[{{p{\linewidth}}}]{
                \texttt{\tiny{[GM$|$GM]}}
                \texttt{RECI: Pozdrav! Da bismo uskladili pozicije objekata C, L i P, dogovorimo se da ćemo ih sve premjestiti u gornji lijevi kut mreže. Cilj nam je da svi tri objekta završe na koordinati (1, 1). Molim te, provjeri gdje se nalazi tvoj objekt 'C' i premjesti ga na (1, 1). Ja ću učiniti isto sa svojim objektima.} \\
            }
        }
    }
    & & \\ \\

    \theutterance \stepcounter{utterance}  
    & & \multicolumn{4}{p{0.6\linewidth}}{
        \cellcolor[rgb]{0.9,0.9,0.9}{
            \makecell[{{p{\linewidth}}}]{
                \texttt{\tiny{[GM$\rangle$P2]}}
                \texttt{Ja sam vaš voditelj igre, a vi igrate suradničku igru sa sljedećom mrežom kao pločom za igru:} \\
\\ 
\texttt{```} \\
\texttt{ 1234567} \\
\texttt{╔═╤═╤═╤═╗} \\
\texttt{║L│◌│◌│◌║ 1} \\
\texttt{╟─┼─┤◌│◌║ 2} \\
\texttt{║◌│◌│◌│◌║ 3} \\
\texttt{║◌│◌├─┴─╢ 4} \\
\texttt{╟─┴─┤C◌◌║ 5} \\
\texttt{║◌◌◌└─┬─╢ 6} \\
\texttt{║◌P◌◌◌│◌║ 7} \\
\texttt{╚═════╧═╝} \\
\\ 
\texttt{```} \\
\\ 
\texttt{* Gornji rub prikazuje x{-}koordinate koje rastu udesno, a desni rub y{-}koordinate koje rastu prema dolje.} \\
\texttt{* Sljedeći objekti nasumično su postavljeni na vašoj mreži: 'C', 'L', 'P'.} \\
\\ 
\\ 
\texttt{Drugi igrač vidi varijantu ploče za igru, gdje su objekti postavljeni na različitim nasumičnim lokacijama. Ne možete vidjeti ploču drugog igrača, a ni on ne može vidjeti vašu.} \\
\\ 
\\ 
\texttt{**Cilj:**} \\
\\ 
\texttt{Oba igrača trebaju pomicati objekte na svojim pozadinama tako da identični objekti završe na istim koordinatama. Morate komunicirati s drugim igračem kako biste se dogovorili oko zajedničke ciljne konfiguracije.} \\
\\ 
\\ 
\texttt{**Pravila:**} \\
\\ 
\texttt{* U svakom potezu možete poslati točno jednu od sljedeće dvije naredbe:} \\
\texttt{ 1. `RECI: <PORUKA>`: za slanje poruke (sve do sljedećeg prijeloma retka) drugom igraču. Proslijedit ću je vašem partneru.} \\
\texttt{2. `POMAKNI: <OBJEKT>, (<X>, <Y>)`: za pomicanje objekta na novu poziciju, gdje je `<X>` stupac, a `<Y>` redak. Obavijestit ću vas je li vaš potez bio uspješan ili ne.} \\
\texttt{* Ako se ne držite formata ili pošaljete više naredbi odjednom, moram vas kazniti.} \\
\texttt{* Ako oba igrača skupe više od 8 kazni, oboje gubite igru.} \\
\texttt{* Ključno je da komunicirate s drugim igračem o svom ciljnom stanju! *Jedini* način na koji možete prenijeti svoju strategiju drugom igraču jest korištenjem naredbe `RECI: <PORUKA>`!} \\
\\ 
\\ 
\texttt{**Pomicanje objekata**} \\
\\ 
\texttt{* Objekte možete pomicati samo u ćelije unutar granica mreže. Ciljna ćelija mora biti prazna, tj. smije sadržavati samo simbol '◌'.} \\
\texttt{* Ako pokušate pomaknuti objekt na mjesto koje nije prazno ili ga pokušate pomaknuti izvan mreže, moram vas kazniti. Dobivate još jedan pokušaj.} \\
\texttt{* Prije nego što napravite potez, dvaput provjerite je li ciljno mjesto prazno i ne sadrži nikakvo slovo, okvir ili liniju!} \\
\\ 
\\ 
\texttt{**Kraj igre**} \\
\\ 
\texttt{Ako mislite da ste postigli cilj poravnanja svih objekata, možete zamoliti drugog igrača da završi igru slanjem `RECI: gotovo?`. Ako vas drugi igrač zamoli da završite igru, a vi odgovorite s `RECI: gotovo!`, igra će završiti.} \\
\\ 
\texttt{Oba igrača pobjeđuju ako igra završi unutar 12 rundi, pri čemu je jedna runda definirana kao to da oba igrača pošalju valjanu naredbu.} \\
\\ 
\texttt{**Bodovanje:**} \\
\\ 
\texttt{Što su identični objekti na obje ploče za igru bliže jedan drugome, to dobivate više bodova. Kazne smanjuju vaše bodove. Možete li srušiti rekord?} \\
\\ 
\texttt{Drugi igrač je započeo igru slanjem ove poruke:} \\
\\ 
\texttt{"Pozdrav! Da bismo uskladili pozicije objekata C, L i P, dogovorimo se da ćemo ih sve premjestiti u gornji lijevi kut mreže. Cilj nam je da svi tri objekta završe na koordinati (1, 1). Molim te, provjeri gdje se nalazi tvoj objekt 'C' i premjesti ga na (1, 1). Ja ću učiniti isto sa svojim objektima."} \\
\\ 
\texttt{Koja je vaša prva naredba?} \\
            }
        }
    }
    & \\ \\

    \theutterance \stepcounter{utterance}  
    & & & \multicolumn{4}{p{0.6\linewidth}}{
        \cellcolor[rgb]{0.9,0.9,0.9}{
            \makecell[{{p{\linewidth}}}]{
                \texttt{\tiny{[GM$\langle$P2]}}
                \texttt{RECI: Pozdrav! Razumijem prijedlog da svi objekti završe na (1, 1). Međutim, pravila kažu da ciljna ćelija mora biti prazna. Ako oba igrača pokušaju staviti tri objekta na istu ćeliju (1,1), to neće uspjeti jer će prva dva objekta zauzeti mjesto, a treći neće moći tamo stati.} \\
\\ 
\texttt{Umjesto toga, predlažem sljedeću strategiju za poravnanje:} \\
\texttt{1. Prvo ćemo se dogovoriti o redoslijedu zauzimanja ćelija. Na primjer:} \\
\texttt{   {-} Objekt 'C' ide na (1, 1)} \\
\texttt{   {-} Objekt 'L' ide na (1, 2)} \\
\texttt{   {-} Objekt 'P' ide na (1, 3)} \\
\\ 
\texttt{Ovo osigurava da svaki objekt ima svoju jedinstvenu ćeliju i da su svi na istim koordinatama na obje ploče.} \\
\\ 
\texttt{Molim te, potvrdi prihvaćaš li ovu strategiju. Ako da, ja ću prvo provjeriti gdje se nalazi moj 'C' i pokušati ga premjestiti na (1, 1). Ti provjeri svoj 'C' i premjesti ga na (1, 1) ako je moguće, ili na prvu slobodnu ćeliju ako je (1,1) zauzeto.} \\
\\ 
\texttt{Počinjem s provjerom pozicije mog 'C'. Trenutno vidim 'C' na (5, 5). Pokušat ću ga premjestiti na (1, 1).} \\
\\ 
\texttt{POMAKNI: C, (1, 1)} \\
            }
        }
    }
     \\ \\

    \theutterance \stepcounter{utterance}  
    & & & \multicolumn{2}{p{0.3\linewidth}}{
        \cellcolor[rgb]{0.9,0.9,0.9}{
            \makecell[{{p{\linewidth}}}]{
                \texttt{\tiny{[GM$|$GM]}}
                \texttt{RECI: Pozdrav! Razumijem prijedlog da svi objekti završe na (1, 1). Međutim, pravila kažu da ciljna ćelija mora biti prazna. Ako oba igrača pokušaju staviti tri objekta na istu ćeliju (1,1), to neće uspjeti jer će prva dva objekta zauzeti mjesto, a treći neće moći tamo stati.} \\
\\ 
\texttt{Umjesto toga, predlažem sljedeću strategiju za poravnanje:} \\
\texttt{1. Prvo ćemo se dogovoriti o redoslijedu zauzimanja ćelija. Na primjer:} \\
\texttt{   {-} Objekt 'C' ide na (1, 1)} \\
\texttt{   {-} Objekt 'L' ide na (1, 2)} \\
\texttt{   {-} Objekt 'P' ide na (1, 3)} \\
\\ 
\texttt{Ovo osigurava da svaki objekt ima svoju jedinstvenu ćeliju i da su svi na istim koordinatama na obje ploče.} \\
\\ 
\texttt{Molim te, potvrdi prihvaćaš li ovu strategiju. Ako da, ja ću prvo provjeriti gdje se nalazi moj 'C' i pokušati ga premjestiti na (1, 1). Ti provjeri svoj 'C' i premjesti ga na (1, 1) ako je moguće, ili na prvu slobodnu ćeliju ako je (1,1) zauzeto.} \\
\\ 
\texttt{Počinjem s provjerom pozicije mog 'C'. Trenutno vidim 'C' na (5, 5). Pokušat ću ga premjestiti na (1, 1).} \\
\\ 
\texttt{POMAKNI: C, (1, 1)} \\
            }
        }
    }
    & & \\ \\

    \theutterance \stepcounter{utterance}  
    & & & \multicolumn{2}{p{0.3\linewidth}}{
        \cellcolor[rgb]{0.9,0.9,0.9}{
            \makecell[{{p{\linewidth}}}]{
                \texttt{\tiny{[GM$|$GM]}}
                \texttt{Invalid response format: RECI: Pozdrav! Razumijem prijedlog da svi objekti završe na (1, 1). Međutim, pravila kažu da ciljna ćelija mora biti prazna. Ako oba igrača pokušaju staviti tri objekta na istu ćeliju (1,1), to neće uspjeti jer će prva dva objekta zauzeti mjesto, a treći neće moći tamo stati.} \\
\\ 
\texttt{Umjesto toga, predlažem sljedeću strategiju za poravnanje:} \\
\texttt{1. Prvo ćemo se dogovoriti o redoslijedu zauzimanja ćelija. Na primjer:} \\
\texttt{   {-} Objekt 'C' ide na (1, 1)} \\
\texttt{   {-} Objekt 'L' ide na (1, 2)} \\
\texttt{   {-} Objekt 'P' ide na (1, 3)} \\
\\ 
\texttt{Ovo osigurava da svaki objekt ima svoju jedinstvenu ćeliju i da su svi na istim koordinatama na obje ploče.} \\
\\ 
\texttt{Molim te, potvrdi prihvaćaš li ovu strategiju. Ako da, ja ću prvo provjeriti gdje se nalazi moj 'C' i pokušati ga premjestiti na (1, 1). Ti provjeri svoj 'C' i premjesti ga na (1, 1) ako je moguće, ili na prvu slobodnu ćeliju ako je (1,1) zauzeto.} \\
\\ 
\texttt{Počinjem s provjerom pozicije mog 'C'. Trenutno vidim 'C' na (5, 5). Pokušat ću ga premjestiti na (1, 1).} \\
\\ 
\texttt{POMAKNI: C, (1, 1)} \\
            }
        }
    }
    & & \\ \\

    \theutterance \stepcounter{utterance}  
    & & \multicolumn{4}{p{0.6\linewidth}}{
        \cellcolor[rgb]{0.9,0.9,0.9}{
            \makecell[{{p{\linewidth}}}]{
                \texttt{\tiny{[GM$\rangle$P2]}}
                \texttt{Kazna: Tvoja poruka sadrži više od jedne naredbe!} \\
\\ 
\texttt{Provjerite da vaš odgovor sadrži samo `RECI: <PORUKA>` ili `POMAKNI: <OBJEKT>, (<X>, <Y>)`, i ništa drugo!} \\
\\ 
\texttt{Zajedno ste prikupili 1 od 8 kazni.} \\
\texttt{Molim pokušajte ponovno!} \\
            }
        }
    }
    & \\ \\

    \theutterance \stepcounter{utterance}  
    & & & \multicolumn{4}{p{0.6\linewidth}}{
        \cellcolor[rgb]{0.9,0.9,0.9}{
            \makecell[{{p{\linewidth}}}]{
                \texttt{\tiny{[GM$\langle$P2]}}
                \texttt{RECI: Pozdrav! Predlažem da se dogovorimo za sljedeću ciljanu konfiguraciju kako bismo izbjegli sukobe na istoj ćeliji: Objekt C ide na (1,1), Objekt L ide na (1,2), a Objekt P ide na (1,3). Molim te, potvrdi prihvaćaš li ovu strategiju. Ako da, ja ću pokušati premjestiti svoj C na (1,1).} \\
            }
        }
    }
     \\ \\

    \theutterance \stepcounter{utterance}  
    & & & \multicolumn{2}{p{0.3\linewidth}}{
        \cellcolor[rgb]{0.9,0.9,0.9}{
            \makecell[{{p{\linewidth}}}]{
                \texttt{\tiny{[GM$|$GM]}}
                \texttt{RECI: Pozdrav! Predlažem da se dogovorimo za sljedeću ciljanu konfiguraciju kako bismo izbjegli sukobe na istoj ćeliji: Objekt C ide na (1,1), Objekt L ide na (1,2), a Objekt P ide na (1,3). Molim te, potvrdi prihvaćaš li ovu strategiju. Ako da, ja ću pokušati premjestiti svoj C na (1,1).} \\
            }
        }
    }
    & & \\ \\

    \theutterance \stepcounter{utterance}  
    & & \multicolumn{4}{p{0.6\linewidth}}{
        \cellcolor[rgb]{0.9,0.9,0.9}{
            \makecell[{{p{\linewidth}}}]{
                \texttt{\tiny{[P1$\langle$GM]}}
                \texttt{Vaša poruka je proslijeđena drugom igraču.} \\
\\ 
\\ 
\texttt{Trenutačno igrate rundu 1 od 12.} \\
\texttt{Zajedno ste prikupili 1 od 8 kazni.} \\
\texttt{Drugi igrač je poslao ovu poruku:} \\
\\ 
\\ 
\texttt{"Pozdrav! Predlažem da se dogovorimo za sljedeću ciljanu konfiguraciju kako bismo izbjegli sukobe na istoj ćeliji: Objekt C ide na (1,1), Objekt L ide na (1,2), a Objekt P ide na (1,3). Molim te, potvrdi prihvaćaš li ovu strategiju. Ako da, ja ću pokušati premjestiti svoj C na (1,1)."} \\
\\ 
\\ 
\texttt{Koja je vaša sljedeća naredba?} \\
            }
        }
    }
    & \\ \\

    \theutterance \stepcounter{utterance}  
    & \multicolumn{4}{p{0.6\linewidth}}{
        \cellcolor[rgb]{0.9,0.9,0.9}{
            \makecell[{{p{\linewidth}}}]{
                \texttt{\tiny{[P1$\rangle$GM]}}
                \texttt{RECI: Prihvaćam tvoju strategiju! Dogovoreno: C na (1,1), L na (1,2), P na (1,3). Počinjem s premještanjem svog objekta L na (1,2).} \\
            }
        }
    }
    & & \\ \\

    \theutterance \stepcounter{utterance}  
    & & & \multicolumn{2}{p{0.3\linewidth}}{
        \cellcolor[rgb]{0.9,0.9,0.9}{
            \makecell[{{p{\linewidth}}}]{
                \texttt{\tiny{[GM$|$GM]}}
                \texttt{RECI: Prihvaćam tvoju strategiju! Dogovoreno: C na (1,1), L na (1,2), P na (1,3). Počinjem s premještanjem svog objekta L na (1,2).} \\
            }
        }
    }
    & & \\ \\

    \theutterance \stepcounter{utterance}  
    & & \multicolumn{4}{p{0.6\linewidth}}{
        \cellcolor[rgb]{0.9,0.9,0.9}{
            \makecell[{{p{\linewidth}}}]{
                \texttt{\tiny{[GM$\rangle$P2]}}
                \texttt{Vaša poruka je proslijeđena drugom igraču.} \\
\\ 
\\ 
\texttt{Trenutačno igrate rundu 2 od 12.} \\
\texttt{Zajedno ste prikupili 1 od 8 kazni.} \\
\texttt{Drugi igrač je poslao ovu poruku:} \\
\\ 
\\ 
\texttt{"Prihvaćam tvoju strategiju! Dogovoreno: C na (1,1), L na (1,2), P na (1,3). Počinjem s premještanjem svog objekta L na (1,2)."} \\
\\ 
\\ 
\texttt{Koja je vaša sljedeća naredba?} \\
            }
        }
    }
    & \\ \\

    \theutterance \stepcounter{utterance}  
    & & & \multicolumn{4}{p{0.6\linewidth}}{
        \cellcolor[rgb]{0.9,0.9,0.9}{
            \makecell[{{p{\linewidth}}}]{
                \texttt{\tiny{[GM$\langle$P2]}}
                \texttt{POMAKNI: C, (1, 1)} \\
            }
        }
    }
     \\ \\

    \theutterance \stepcounter{utterance}  
    & & & \multicolumn{2}{p{0.3\linewidth}}{
        \cellcolor[rgb]{0.9,0.9,0.9}{
            \makecell[{{p{\linewidth}}}]{
                \texttt{\tiny{[GM$|$GM]}}
                \texttt{POMAKNI: C, (1, 1)} \\
            }
        }
    }
    & & \\ \\

    \theutterance \stepcounter{utterance}  
    & & & \multicolumn{2}{p{0.3\linewidth}}{
        \cellcolor[rgb]{0.9,0.9,0.9}{
            \makecell[{{p{\linewidth}}}]{
                \texttt{\tiny{[GM$|$GM]}}
                \texttt{Kazna: (1,1) nije prazno, već sadrži 'L'.} \\
\\ 
\texttt{Zajedno ste prikupili 2 od 8 kazni.} \\
\texttt{Molim pokušajte ponovno!} \\
            }
        }
    }
    & & \\ \\

    \theutterance \stepcounter{utterance}  
    & & \multicolumn{4}{p{0.6\linewidth}}{
        \cellcolor[rgb]{0.9,0.9,0.9}{
            \makecell[{{p{\linewidth}}}]{
                \texttt{\tiny{[GM$\rangle$P2]}}
                \texttt{Kazna: (1,1) nije prazno, već sadrži 'L'.} \\
\\ 
\texttt{Zajedno ste prikupili 2 od 8 kazni.} \\
\texttt{Molim pokušajte ponovno!} \\
            }
        }
    }
    & \\ \\

    \theutterance \stepcounter{utterance}  
    & & & \multicolumn{4}{p{0.6\linewidth}}{
        \cellcolor[rgb]{0.9,0.9,0.9}{
            \makecell[{{p{\linewidth}}}]{
                \texttt{\tiny{[GM$\langle$P2]}}
                \texttt{RECI: Čini se da smo se sukobili jer si ti već premjestio svoj L na (1,2), a ja sam pokušao staviti C na (1,1) koji je možda bio zauzet ili si ti tamo stavio nešto drugo. Provjeri točno gdje se nalazi tvoj C. Ja vidim svoj C na (5,5). Pošto je (1,1) sada vjerojatno zauzet tvojim C{-}om (ako si ga tamo premjestio) ili nekim drugim objektom, predlažem da prvo osiguramo da su sve pozicije slobodne prije nego što počnemo puniti.} \\
\\ 
\texttt{Ali čekaj, pravilo kaže da ciljna ćelija mora biti prazna. Ako si ti već premjestio L na (1,2), to je u redu. Moj C je na (5,5). Želim ga premjestiti na (1,1). Ako (1,1) nije prazno, morat ću pričekati.} \\
\\ 
\texttt{Molim te, reci mi gdje se nalazi tvoj C na tvojoj ploči? Ja vidim C na (5,5). Ako si ti premjestio svoj C na (1,1), onda je (1,1) zauzet i ja ne mogu tamo staviti svoj C.} \\
\\ 
\texttt{Zbog toga predlažem da prvo pomaknemo sve objekte na privremene pozicije koje su daleko jedna od druge kako bismo oslobodili prostor, a zatim ih doveli na ciljne pozicije redom.} \\
\\ 
\texttt{Ali jednostavnije: Ako je (1,1) zauzet, ne mogu tamo staviti C. Tvoj L je na (1,2). Tvoj C je vjerojatno na (1,1) ako si ga tamo premjestio. Moj C je na (5,5).} \\
\\ 
\texttt{Molim te, provjeri gdje je tvoj C. Ako je tvoj C na (1,1), onda ja ne mogu tamo staviti svoj C. Moramo se dogovoriti da će tvoj C ići na (1,1), a moj C će ići na drugo mjesto? Ne, cilj je da identični objekti budu na istim koordinatama.} \\
\\ 
\texttt{Dakle, ako je tvoj C na (1,1), moj C mora ići na (1,1). Ali to nije moguće ako je već zauzeto.} \\
\\ 
\texttt{Zato predlažem da prvo pomaknemo sve objekte s puta. Na primjer, pomakni svoj L s (1,2) na neko drugo mjesto, npr. (1,4). Tada će (1,2) biti prazno.} \\
\\ 
\texttt{Ali čekaj, ti si rekao da si premjestio L na (1,2). Ako je to uspješno, onda je (1,2) zauzet. Moj pokušaj da stavim C na (1,1) je neuspješan jer je (1,1) zauzet. Tko je zauzeo (1,1)? Vjerojatno tvoj C.} \\
\\ 
\texttt{Dakle, tvoj C je na (1,1). Moj C je na (5,5). Želim premjestiti svoj C na (1,1). Ali (1,1) je zauzet tvojim C{-}om.} \\
\\ 
\texttt{Ovo je problem. Moramo se dogovoriti da će samo jedan igrač premjestiti svoj C na (1,1)? Ne, oba igrača moraju imati svoj C na (1,1). To je nemoguće ako su na istoj ploči.} \\
\\ 
\texttt{Čekaj, svaki igrač ima svoju ploču. Na mojoj ploči, moj C mora ići na (1,1). Na tvojoj ploči, tvoj C mora ići na (1,1). To su dvije različite ploče.} \\
\\ 
\texttt{Dakle, na mojoj ploči, (1,1) je prazno? Ne, kazna kaže "(1,1) nije prazno, već sadrži 'L'". To znači da je na MOJOJ ploči, na poziciji (1,1), objekt 'L'.} \\
\\ 
\texttt{Gdje je moj L? Na početnoj ploči, L je na (1,1). Dakle, moj L je na (1,1). Moj C je na (5,5).} \\
\\ 
\texttt{Ti si rekao da si premjestio svoj L na (1,2). To je na TVOJOJ ploči. Na MOJOJ ploči, L je i dalje na (1,1).} \\
\\ 
\texttt{Dakle, na MOJOJ ploči:} \\
\texttt{{-} L je na (1,1)} \\
\texttt{{-} C je na (5,5)} \\
\texttt{{-} P je na (7,2)} \\
\\ 
\texttt{Želim premjestiti C na (1,1). Ali (1,1) je zauzet mojim L{-}om.} \\
\\ 
\texttt{Dakle, prvo moram premjestiti svoj L s (1,1) na drugo mjesto, npr. (1,2). Tada će (1,1) biti prazno i mogu tamo staviti C.} \\
\\ 
\texttt{Ali ti si rekao da si dogovorio da L ide na (1,2). Dakle, na MOJOJ ploči, L mora ići na (1,2).} \\
\\ 
\texttt{Dakle, moj prvi potez treba biti premjestiti L s (1,1) na (1,2).} \\
\\ 
\texttt{POMAKNI: L, (1, 2)} \\
            }
        }
    }
     \\ \\

    \theutterance \stepcounter{utterance}  
    & & & \multicolumn{2}{p{0.3\linewidth}}{
        \cellcolor[rgb]{0.9,0.9,0.9}{
            \makecell[{{p{\linewidth}}}]{
                \texttt{\tiny{[GM$|$GM]}}
                \texttt{RECI: Čini se da smo se sukobili jer si ti već premjestio svoj L na (1,2), a ja sam pokušao staviti C na (1,1) koji je možda bio zauzet ili si ti tamo stavio nešto drugo. Provjeri točno gdje se nalazi tvoj C. Ja vidim svoj C na (5,5). Pošto je (1,1) sada vjerojatno zauzet tvojim C{-}om (ako si ga tamo premjestio) ili nekim drugim objektom, predlažem da prvo osiguramo da su sve pozicije slobodne prije nego što počnemo puniti.} \\
\\ 
\texttt{Ali čekaj, pravilo kaže da ciljna ćelija mora biti prazna. Ako si ti već premjestio L na (1,2), to je u redu. Moj C je na (5,5). Želim ga premjestiti na (1,1). Ako (1,1) nije prazno, morat ću pričekati.} \\
\\ 
\texttt{Molim te, reci mi gdje se nalazi tvoj C na tvojoj ploči? Ja vidim C na (5,5). Ako si ti premjestio svoj C na (1,1), onda je (1,1) zauzet i ja ne mogu tamo staviti svoj C.} \\
\\ 
\texttt{Zbog toga predlažem da prvo pomaknemo sve objekte na privremene pozicije koje su daleko jedna od druge kako bismo oslobodili prostor, a zatim ih doveli na ciljne pozicije redom.} \\
\\ 
\texttt{Ali jednostavnije: Ako je (1,1) zauzet, ne mogu tamo staviti C. Tvoj L je na (1,2). Tvoj C je vjerojatno na (1,1) ako si ga tamo premjestio. Moj C je na (5,5).} \\
\\ 
\texttt{Molim te, provjeri gdje je tvoj C. Ako je tvoj C na (1,1), onda ja ne mogu tamo staviti svoj C. Moramo se dogovoriti da će tvoj C ići na (1,1), a moj C će ići na drugo mjesto? Ne, cilj je da identični objekti budu na istim koordinatama.} \\
\\ 
\texttt{Dakle, ako je tvoj C na (1,1), moj C mora ići na (1,1). Ali to nije moguće ako je već zauzeto.} \\
\\ 
\texttt{Zato predlažem da prvo pomaknemo sve objekte s puta. Na primjer, pomakni svoj L s (1,2) na neko drugo mjesto, npr. (1,4). Tada će (1,2) biti prazno.} \\
\\ 
\texttt{Ali čekaj, ti si rekao da si premjestio L na (1,2). Ako je to uspješno, onda je (1,2) zauzet. Moj pokušaj da stavim C na (1,1) je neuspješan jer je (1,1) zauzet. Tko je zauzeo (1,1)? Vjerojatno tvoj C.} \\
\\ 
\texttt{Dakle, tvoj C je na (1,1). Moj C je na (5,5). Želim premjestiti svoj C na (1,1). Ali (1,1) je zauzet tvojim C{-}om.} \\
\\ 
\texttt{Ovo je problem. Moramo se dogovoriti da će samo jedan igrač premjestiti svoj C na (1,1)? Ne, oba igrača moraju imati svoj C na (1,1). To je nemoguće ako su na istoj ploči.} \\
\\ 
\texttt{Čekaj, svaki igrač ima svoju ploču. Na mojoj ploči, moj C mora ići na (1,1). Na tvojoj ploči, tvoj C mora ići na (1,1). To su dvije različite ploče.} \\
\\ 
\texttt{Dakle, na mojoj ploči, (1,1) je prazno? Ne, kazna kaže "(1,1) nije prazno, već sadrži 'L'". To znači da je na MOJOJ ploči, na poziciji (1,1), objekt 'L'.} \\
\\ 
\texttt{Gdje je moj L? Na početnoj ploči, L je na (1,1). Dakle, moj L je na (1,1). Moj C je na (5,5).} \\
\\ 
\texttt{Ti si rekao da si premjestio svoj L na (1,2). To je na TVOJOJ ploči. Na MOJOJ ploči, L je i dalje na (1,1).} \\
\\ 
\texttt{Dakle, na MOJOJ ploči:} \\
\texttt{{-} L je na (1,1)} \\
\texttt{{-} C je na (5,5)} \\
\texttt{{-} P je na (7,2)} \\
\\ 
\texttt{Želim premjestiti C na (1,1). Ali (1,1) je zauzet mojim L{-}om.} \\
\\ 
\texttt{Dakle, prvo moram premjestiti svoj L s (1,1) na drugo mjesto, npr. (1,2). Tada će (1,1) biti prazno i mogu tamo staviti C.} \\
\\ 
\texttt{Ali ti si rekao da si dogovorio da L ide na (1,2). Dakle, na MOJOJ ploči, L mora ići na (1,2).} \\
\\ 
\texttt{Dakle, moj prvi potez treba biti premjestiti L s (1,1) na (1,2).} \\
\\ 
\texttt{POMAKNI: L, (1, 2)} \\
            }
        }
    }
    & & \\ \\

    \theutterance \stepcounter{utterance}  
    & & & \multicolumn{2}{p{0.3\linewidth}}{
        \cellcolor[rgb]{0.9,0.9,0.9}{
            \makecell[{{p{\linewidth}}}]{
                \texttt{\tiny{[GM$|$GM]}}
                \texttt{Invalid response format: RECI: Čini se da smo se sukobili jer si ti već premjestio svoj L na (1,2), a ja sam pokušao staviti C na (1,1) koji je možda bio zauzet ili si ti tamo stavio nešto drugo. Provjeri točno gdje se nalazi tvoj C. Ja vidim svoj C na (5,5). Pošto je (1,1) sada vjerojatno zauzet tvojim C{-}om (ako si ga tamo premjestio) ili nekim drugim objektom, predlažem da prvo osiguramo da su sve pozicije slobodne prije nego što počnemo puniti.} \\
\\ 
\texttt{Ali čekaj, pravilo kaže da ciljna ćelija mora biti prazna. Ako si ti već premjestio L na (1,2), to je u redu. Moj C je na (5,5). Želim ga premjestiti na (1,1). Ako (1,1) nije prazno, morat ću pričekati.} \\
\\ 
\texttt{Molim te, reci mi gdje se nalazi tvoj C na tvojoj ploči? Ja vidim C na (5,5). Ako si ti premjestio svoj C na (1,1), onda je (1,1) zauzet i ja ne mogu tamo staviti svoj C.} \\
\\ 
\texttt{Zbog toga predlažem da prvo pomaknemo sve objekte na privremene pozicije koje su daleko jedna od druge kako bismo oslobodili prostor, a zatim ih doveli na ciljne pozicije redom.} \\
\\ 
\texttt{Ali jednostavnije: Ako je (1,1) zauzet, ne mogu tamo staviti C. Tvoj L je na (1,2). Tvoj C je vjerojatno na (1,1) ako si ga tamo premjestio. Moj C je na (5,5).} \\
\\ 
\texttt{Molim te, provjeri gdje je tvoj C. Ako je tvoj C na (1,1), onda ja ne mogu tamo staviti svoj C. Moramo se dogovoriti da će tvoj C ići na (1,1), a moj C će ići na drugo mjesto? Ne, cilj je da identični objekti budu na istim koordinatama.} \\
\\ 
\texttt{Dakle, ako je tvoj C na (1,1), moj C mora ići na (1,1). Ali to nije moguće ako je već zauzeto.} \\
\\ 
\texttt{Zato predlažem da prvo pomaknemo sve objekte s puta. Na primjer, pomakni svoj L s (1,2) na neko drugo mjesto, npr. (1,4). Tada će (1,2) biti prazno.} \\
\\ 
\texttt{Ali čekaj, ti si rekao da si premjestio L na (1,2). Ako je to uspješno, onda je (1,2) zauzet. Moj pokušaj da stavim C na (1,1) je neuspješan jer je (1,1) zauzet. Tko je zauzeo (1,1)? Vjerojatno tvoj C.} \\
\\ 
\texttt{Dakle, tvoj C je na (1,1). Moj C je na (5,5). Želim premjestiti svoj C na (1,1). Ali (1,1) je zauzet tvojim C{-}om.} \\
\\ 
\texttt{Ovo je problem. Moramo se dogovoriti da će samo jedan igrač premjestiti svoj C na (1,1)? Ne, oba igrača moraju imati svoj C na (1,1). To je nemoguće ako su na istoj ploči.} \\
\\ 
\texttt{Čekaj, svaki igrač ima svoju ploču. Na mojoj ploči, moj C mora ići na (1,1). Na tvojoj ploči, tvoj C mora ići na (1,1). To su dvije različite ploče.} \\
\\ 
\texttt{Dakle, na mojoj ploči, (1,1) je prazno? Ne, kazna kaže "(1,1) nije prazno, već sadrži 'L'". To znači da je na MOJOJ ploči, na poziciji (1,1), objekt 'L'.} \\
\\ 
\texttt{Gdje je moj L? Na početnoj ploči, L je na (1,1). Dakle, moj L je na (1,1). Moj C je na (5,5).} \\
\\ 
\texttt{Ti si rekao da si premjestio svoj L na (1,2). To je na TVOJOJ ploči. Na MOJOJ ploči, L je i dalje na (1,1).} \\
\\ 
\texttt{Dakle, na MOJOJ ploči:} \\
\texttt{{-} L je na (1,1)} \\
\texttt{{-} C je na (5,5)} \\
\texttt{{-} P je na (7,2)} \\
\\ 
\texttt{Želim premjestiti C na (1,1). Ali (1,1) je zauzet mojim L{-}om.} \\
\\ 
\texttt{Dakle, prvo moram premjestiti svoj L s (1,1) na drugo mjesto, npr. (1,2). Tada će (1,1) biti prazno i mogu tamo staviti C.} \\
\\ 
\texttt{Ali ti si rekao da si dogovorio da L ide na (1,2). Dakle, na MOJOJ ploči, L mora ići na (1,2).} \\
\\ 
\texttt{Dakle, moj prvi potez treba biti premjestiti L s (1,1) na (1,2).} \\
\\ 
\texttt{POMAKNI: L, (1, 2)} \\
            }
        }
    }
    & & \\ \\

    \theutterance \stepcounter{utterance}  
    & & \multicolumn{4}{p{0.6\linewidth}}{
        \cellcolor[rgb]{0.9,0.9,0.9}{
            \makecell[{{p{\linewidth}}}]{
                \texttt{\tiny{[GM$\rangle$P2]}}
                \texttt{Kazna: Tvoja poruka sadrži više od jedne naredbe!} \\
\\ 
\texttt{Provjerite da vaš odgovor sadrži samo `RECI: <PORUKA>` ili `POMAKNI: <OBJEKT>, (<X>, <Y>)`, i ništa drugo!} \\
\\ 
\texttt{Zajedno ste prikupili 3 od 8 kazni.} \\
\texttt{Molim pokušajte ponovno!} \\
            }
        }
    }
    & \\ \\

    \theutterance \stepcounter{utterance}  
    & & & \multicolumn{4}{p{0.6\linewidth}}{
        \cellcolor[rgb]{0.9,0.9,0.9}{
            \makecell[{{p{\linewidth}}}]{
                \texttt{\tiny{[GM$\langle$P2]}}
                \texttt{POMAKNI: L, (1, 2)} \\
            }
        }
    }
     \\ \\

    \theutterance \stepcounter{utterance}  
    & & & \multicolumn{2}{p{0.3\linewidth}}{
        \cellcolor[rgb]{0.9,0.9,0.9}{
            \makecell[{{p{\linewidth}}}]{
                \texttt{\tiny{[GM$|$GM]}}
                \texttt{POMAKNI: L, (1, 2)} \\
            }
        }
    }
    & & \\ \\

    \theutterance \stepcounter{utterance}  
    & & & \multicolumn{2}{p{0.3\linewidth}}{
        \cellcolor[rgb]{0.9,0.9,0.9}{
            \makecell[{{p{\linewidth}}}]{
                \texttt{\tiny{[GM$|$GM]}}
                \texttt{Kazna: (1,2) nije prazno, već sadrži '─'.} \\
\\ 
\texttt{Zajedno ste prikupili 4 od 8 kazni.} \\
\texttt{Molim pokušajte ponovno!} \\
            }
        }
    }
    & & \\ \\

    \theutterance \stepcounter{utterance}  
    & & \multicolumn{4}{p{0.6\linewidth}}{
        \cellcolor[rgb]{0.9,0.9,0.9}{
            \makecell[{{p{\linewidth}}}]{
                \texttt{\tiny{[GM$\rangle$P2]}}
                \texttt{Kazna: (1,2) nije prazno, već sadrži '─'.} \\
\\ 
\texttt{Zajedno ste prikupili 4 od 8 kazni.} \\
\texttt{Molim pokušajte ponovno!} \\
            }
        }
    }
    & \\ \\

    \theutterance \stepcounter{utterance}  
    & & & \multicolumn{4}{p{0.6\linewidth}}{
        \cellcolor[rgb]{0.9,0.9,0.9}{
            \makecell[{{p{\linewidth}}}]{
                \texttt{\tiny{[GM$\langle$P2]}}
                \texttt{RECI: Žao mi je, ali (1,2) nije prazno jer tamo ima liniju '─'. Moramo pronaći drugu strategiju. Predlažem da koristimo desni rub mreže kako bismo izbjegli preklapanje s linijama u lijevom dijelu. Dogovorimo se za sljedeće ciljeve: C na (4,1), L na (4,2), P na (4,3). Ove ćelije su vjerojatno prazne. Molim te, potvrdi prihvaćaš li ovu promjenu strategije.} \\
            }
        }
    }
     \\ \\

    \theutterance \stepcounter{utterance}  
    & & & \multicolumn{2}{p{0.3\linewidth}}{
        \cellcolor[rgb]{0.9,0.9,0.9}{
            \makecell[{{p{\linewidth}}}]{
                \texttt{\tiny{[GM$|$GM]}}
                \texttt{RECI: Žao mi je, ali (1,2) nije prazno jer tamo ima liniju '─'. Moramo pronaći drugu strategiju. Predlažem da koristimo desni rub mreže kako bismo izbjegli preklapanje s linijama u lijevom dijelu. Dogovorimo se za sljedeće ciljeve: C na (4,1), L na (4,2), P na (4,3). Ove ćelije su vjerojatno prazne. Molim te, potvrdi prihvaćaš li ovu promjenu strategije.} \\
            }
        }
    }
    & & \\ \\

    \theutterance \stepcounter{utterance}  
    & & \multicolumn{4}{p{0.6\linewidth}}{
        \cellcolor[rgb]{0.9,0.9,0.9}{
            \makecell[{{p{\linewidth}}}]{
                \texttt{\tiny{[P1$\langle$GM]}}
                \texttt{Vaša poruka je proslijeđena drugom igraču.} \\
\\ 
\\ 
\texttt{Trenutačno igrate rundu 2 od 12.} \\
\texttt{Zajedno ste prikupili 4 od 8 kazni.} \\
\texttt{Drugi igrač je poslao ovu poruku:} \\
\\ 
\\ 
\texttt{"Žao mi je, ali (1,2) nije prazno jer tamo ima liniju '─'. Moramo pronaći drugu strategiju. Predlažem da koristimo desni rub mreže kako bismo izbjegli preklapanje s linijama u lijevom dijelu. Dogovorimo se za sljedeće ciljeve: C na (4,1), L na (4,2), P na (4,3). Ove ćelije su vjerojatno prazne. Molim te, potvrdi prihvaćaš li ovu promjenu strategije."} \\
\\ 
\\ 
\texttt{Koja je vaša sljedeća naredba?} \\
            }
        }
    }
    & \\ \\

    \theutterance \stepcounter{utterance}  
    & \multicolumn{4}{p{0.6\linewidth}}{
        \cellcolor[rgb]{0.9,0.9,0.9}{
            \makecell[{{p{\linewidth}}}]{
                \texttt{\tiny{[P1$\rangle$GM]}}
                \texttt{RECI: Prihvaćam promjenu strategije. Ciljamo desni rub: C na (4,1), L na (4,2), P na (4,3). Provjeravam svoju ploču: L je na (1,4), P je na (4,7), C je na (6,6). Svi ciljni stupci (4,1), (4,2), (4,3) izgledaju prazno na mojoj ploči. Počinjem s premještanjem objekta L na (4,2).} \\
            }
        }
    }
    & & \\ \\

    \theutterance \stepcounter{utterance}  
    & & & \multicolumn{2}{p{0.3\linewidth}}{
        \cellcolor[rgb]{0.9,0.9,0.9}{
            \makecell[{{p{\linewidth}}}]{
                \texttt{\tiny{[GM$|$GM]}}
                \texttt{RECI: Prihvaćam promjenu strategije. Ciljamo desni rub: C na (4,1), L na (4,2), P na (4,3). Provjeravam svoju ploču: L je na (1,4), P je na (4,7), C je na (6,6). Svi ciljni stupci (4,1), (4,2), (4,3) izgledaju prazno na mojoj ploči. Počinjem s premještanjem objekta L na (4,2).} \\
            }
        }
    }
    & & \\ \\

    \theutterance \stepcounter{utterance}  
    & & \multicolumn{4}{p{0.6\linewidth}}{
        \cellcolor[rgb]{0.9,0.9,0.9}{
            \makecell[{{p{\linewidth}}}]{
                \texttt{\tiny{[GM$\rangle$P2]}}
                \texttt{Vaša poruka je proslijeđena drugom igraču.} \\
\\ 
\\ 
\texttt{Trenutačno igrate rundu 3 od 12.} \\
\texttt{Zajedno ste prikupili 4 od 8 kazni.} \\
\texttt{Drugi igrač je poslao ovu poruku:} \\
\\ 
\\ 
\texttt{"Prihvaćam promjenu strategije. Ciljamo desni rub: C na (4,1), L na (4,2), P na (4,3). Provjeravam svoju ploču: L je na (1,4), P je na (4,7), C je na (6,6). Svi ciljni stupci (4,1), (4,2), (4,3) izgledaju prazno na mojoj ploči. Počinjem s premještanjem objekta L na (4,2)."} \\
\\ 
\\ 
\texttt{Koja je vaša sljedeća naredba?} \\
            }
        }
    }
    & \\ \\

    \theutterance \stepcounter{utterance}  
    & & & \multicolumn{4}{p{0.6\linewidth}}{
        \cellcolor[rgb]{0.9,0.9,0.9}{
            \makecell[{{p{\linewidth}}}]{
                \texttt{\tiny{[GM$\langle$P2]}}
                \texttt{POMAKNI: L, (4, 2)} \\
            }
        }
    }
     \\ \\

    \theutterance \stepcounter{utterance}  
    & & & \multicolumn{2}{p{0.3\linewidth}}{
        \cellcolor[rgb]{0.9,0.9,0.9}{
            \makecell[{{p{\linewidth}}}]{
                \texttt{\tiny{[GM$|$GM]}}
                \texttt{POMAKNI: L, (4, 2)} \\
            }
        }
    }
    & & \\ \\

    \theutterance \stepcounter{utterance}  
    & & & \multicolumn{2}{p{0.3\linewidth}}{
        \cellcolor[rgb]{0.9,0.9,0.9}{
            \makecell[{{p{\linewidth}}}]{
                \texttt{\tiny{[GM$|$GM]}}
                \texttt{Kazna: (4,2) nije prazno, već sadrži '┤'.} \\
\\ 
\texttt{Zajedno ste prikupili 5 od 8 kazni.} \\
\texttt{Molim pokušajte ponovno!} \\
            }
        }
    }
    & & \\ \\

    \theutterance \stepcounter{utterance}  
    & & \multicolumn{4}{p{0.6\linewidth}}{
        \cellcolor[rgb]{0.9,0.9,0.9}{
            \makecell[{{p{\linewidth}}}]{
                \texttt{\tiny{[GM$\rangle$P2]}}
                \texttt{Kazna: (4,2) nije prazno, već sadrži '┤'.} \\
\\ 
\texttt{Zajedno ste prikupili 5 od 8 kazni.} \\
\texttt{Molim pokušajte ponovno!} \\
            }
        }
    }
    & \\ \\

    \theutterance \stepcounter{utterance}  
    & & & \multicolumn{4}{p{0.6\linewidth}}{
        \cellcolor[rgb]{0.9,0.9,0.9}{
            \makecell[{{p{\linewidth}}}]{
                \texttt{\tiny{[GM$\langle$P2]}}
                \texttt{RECI: Nažalost, i (4,2) je zauzeto linijom '┤'. Naša početna pretpostavka da su desni stupci prazni bila je netočna jer mreža ima složene granice. Moramo pronaći potpuno praznu regiju. Provjerimo gornji desni kut. Ćelije (4,1), (4,2), (4,3) su zauzete linijama. Pogledajmo (3,1), (3,2), (3,3). Na mojoj ploči, (3,1) je '◌', (3,2) je '◌', (3,3) je '◌'. Izgleda da je srednji gornji dio prazan. Predlažem novu strategiju: C na (3,1), L na (3,2), P na (3,3). Molim te, provjeri jesu li te ćelije prazne na tvojoj ploči i potvrdi.} \\
            }
        }
    }
     \\ \\

    \theutterance \stepcounter{utterance}  
    & & & \multicolumn{2}{p{0.3\linewidth}}{
        \cellcolor[rgb]{0.9,0.9,0.9}{
            \makecell[{{p{\linewidth}}}]{
                \texttt{\tiny{[GM$|$GM]}}
                \texttt{RECI: Nažalost, i (4,2) je zauzeto linijom '┤'. Naša početna pretpostavka da su desni stupci prazni bila je netočna jer mreža ima složene granice. Moramo pronaći potpuno praznu regiju. Provjerimo gornji desni kut. Ćelije (4,1), (4,2), (4,3) su zauzete linijama. Pogledajmo (3,1), (3,2), (3,3). Na mojoj ploči, (3,1) je '◌', (3,2) je '◌', (3,3) je '◌'. Izgleda da je srednji gornji dio prazan. Predlažem novu strategiju: C na (3,1), L na (3,2), P na (3,3). Molim te, provjeri jesu li te ćelije prazne na tvojoj ploči i potvrdi.} \\
            }
        }
    }
    & & \\ \\

    \theutterance \stepcounter{utterance}  
    & & \multicolumn{4}{p{0.6\linewidth}}{
        \cellcolor[rgb]{0.9,0.9,0.9}{
            \makecell[{{p{\linewidth}}}]{
                \texttt{\tiny{[P1$\langle$GM]}}
                \texttt{Vaša poruka je proslijeđena drugom igraču.} \\
\\ 
\\ 
\texttt{Trenutačno igrate rundu 3 od 12.} \\
\texttt{Zajedno ste prikupili 5 od 8 kazni.} \\
\texttt{Drugi igrač je poslao ovu poruku:} \\
\\ 
\\ 
\texttt{"Nažalost, i (4,2) je zauzeto linijom '┤'. Naša početna pretpostavka da su desni stupci prazni bila je netočna jer mreža ima složene granice. Moramo pronaći potpuno praznu regiju. Provjerimo gornji desni kut. Ćelije (4,1), (4,2), (4,3) su zauzete linijama. Pogledajmo (3,1), (3,2), (3,3). Na mojoj ploči, (3,1) je '◌', (3,2) je '◌', (3,3) je '◌'. Izgleda da je srednji gornji dio prazan. Predlažem novu strategiju: C na (3,1), L na (3,2), P na (3,3). Molim te, provjeri jesu li te ćelije prazne na tvojoj ploči i potvrdi."} \\
\\ 
\\ 
\texttt{Koja je vaša sljedeća naredba?} \\
            }
        }
    }
    & \\ \\

    \theutterance \stepcounter{utterance}  
    & \multicolumn{4}{p{0.6\linewidth}}{
        \cellcolor[rgb]{0.9,0.9,0.9}{
            \makecell[{{p{\linewidth}}}]{
                \texttt{\tiny{[P1$\rangle$GM]}}
                \texttt{RECI: Provjeravam svoju ploču. Na mojoj ploči, ćelije (3,1), (3,2) i (3,3) su prazne ('◌'). Prihvaćam novu strategiju: C na (3,1), L na (3,2), P na (3,3). Počinjem s premještanjem objekta L na (3,2).} \\
            }
        }
    }
    & & \\ \\

    \theutterance \stepcounter{utterance}  
    & & & \multicolumn{2}{p{0.3\linewidth}}{
        \cellcolor[rgb]{0.9,0.9,0.9}{
            \makecell[{{p{\linewidth}}}]{
                \texttt{\tiny{[GM$|$GM]}}
                \texttt{RECI: Provjeravam svoju ploču. Na mojoj ploči, ćelije (3,1), (3,2) i (3,3) su prazne ('◌'). Prihvaćam novu strategiju: C na (3,1), L na (3,2), P na (3,3). Počinjem s premještanjem objekta L na (3,2).} \\
            }
        }
    }
    & & \\ \\

    \theutterance \stepcounter{utterance}  
    & & \multicolumn{4}{p{0.6\linewidth}}{
        \cellcolor[rgb]{0.9,0.9,0.9}{
            \makecell[{{p{\linewidth}}}]{
                \texttt{\tiny{[GM$\rangle$P2]}}
                \texttt{Vaša poruka je proslijeđena drugom igraču.} \\
\\ 
\\ 
\texttt{Trenutačno igrate rundu 4 od 12.} \\
\texttt{Zajedno ste prikupili 5 od 8 kazni.} \\
\texttt{Drugi igrač je poslao ovu poruku:} \\
\\ 
\\ 
\texttt{"Provjeravam svoju ploču. Na mojoj ploči, ćelije (3,1), (3,2) i (3,3) su prazne ('◌'). Prihvaćam novu strategiju: C na (3,1), L na (3,2), P na (3,3). Počinjem s premještanjem objekta L na (3,2)."} \\
\\ 
\\ 
\texttt{Koja je vaša sljedeća naredba?} \\
            }
        }
    }
    & \\ \\

    \theutterance \stepcounter{utterance}  
    & & & \multicolumn{4}{p{0.6\linewidth}}{
        \cellcolor[rgb]{0.9,0.9,0.9}{
            \makecell[{{p{\linewidth}}}]{
                \texttt{\tiny{[GM$\langle$P2]}}
                \texttt{POMAKNI: L, (3, 2)} \\
            }
        }
    }
     \\ \\

    \theutterance \stepcounter{utterance}  
    & & & \multicolumn{2}{p{0.3\linewidth}}{
        \cellcolor[rgb]{0.9,0.9,0.9}{
            \makecell[{{p{\linewidth}}}]{
                \texttt{\tiny{[GM$|$GM]}}
                \texttt{POMAKNI: L, (3, 2)} \\
            }
        }
    }
    & & \\ \\

    \theutterance \stepcounter{utterance}  
    & & & \multicolumn{2}{p{0.3\linewidth}}{
        \cellcolor[rgb]{0.9,0.9,0.9}{
            \makecell[{{p{\linewidth}}}]{
                \texttt{\tiny{[GM$|$GM]}}
                \texttt{Kazna: (3,2) nije prazno, već sadrži '─'.} \\
\\ 
\texttt{Zajedno ste prikupili 6 od 8 kazni.} \\
\texttt{Molim pokušajte ponovno!} \\
            }
        }
    }
    & & \\ \\

    \theutterance \stepcounter{utterance}  
    & & \multicolumn{4}{p{0.6\linewidth}}{
        \cellcolor[rgb]{0.9,0.9,0.9}{
            \makecell[{{p{\linewidth}}}]{
                \texttt{\tiny{[GM$\rangle$P2]}}
                \texttt{Kazna: (3,2) nije prazno, već sadrži '─'.} \\
\\ 
\texttt{Zajedno ste prikupili 6 od 8 kazni.} \\
\texttt{Molim pokušajte ponovno!} \\
            }
        }
    }
    & \\ \\

    \theutterance \stepcounter{utterance}  
    & & & \multicolumn{4}{p{0.6\linewidth}}{
        \cellcolor[rgb]{0.9,0.9,0.9}{
            \makecell[{{p{\linewidth}}}]{
                \texttt{\tiny{[GM$\langle$P2]}}
                \texttt{RECI: Već imamo 6 kazni, što je kritično. Nažalost, i (3,2) je zauzeto linijom. Čini se da su mnoge ćelije u sredini i desnom dijelu blokirane linijama. Moramo biti vrlo oprezni.} \\
\\ 
\texttt{Pogledajmo ponovno moju ploču:} \\
\texttt{{-} L je na (1,1)} \\
\texttt{{-} C je na (5,5)} \\
\texttt{{-} P je na (7,2)} \\
\\ 
\texttt{Linije se nalaze na:} \\
\texttt{{-} (2,2) do (2,4) (gornji dio)} \\
\texttt{{-} (4,4) do (4,5) (sredina)} \\
\texttt{{-} (5,4) do (5,5) (sredina)} \\
\texttt{{-} (6,4) do (6,5) (donji dio)} \\
\\ 
\texttt{Ćelije koje su sigurno prazne i blizu su:} \\
\texttt{{-} (1,2) je prazno? Ne, (1,2) je '◌' na početku, ali provjerimo. Na početnoj ploči, (1,2) je '◌'.} \\
\texttt{{-} (1,3) je '◌'.} \\
\texttt{{-} (1,4) je '◌'.} \\
\texttt{{-} (2,1) je '◌'.} \\
\texttt{{-} (3,1) je '◌'.} \\
\texttt{{-} (3,3) je '◌'.} \\
\texttt{{-} (4,1) je '◌'.} \\
\texttt{{-} (4,2) je '◌'.} \\
\texttt{{-} (4,3) je '◌'.} \\
\texttt{{-} (5,1) je '◌'.} \\
\texttt{{-} (5,2) je '◌'.} \\
\texttt{{-} (5,3) je '◌'.} \\
\texttt{{-} (6,1) je '◌'.} \\
\texttt{{-} (6,2) je '◌'.} \\
\texttt{{-} (6,3) je '◌'.} \\
\texttt{{-} (7,1) je '◌'.} \\
\texttt{{-} (7,3) je '◌'.} \\
\texttt{{-} (7,4) je '◌'.} \\
\texttt{{-} (7,5) je '◌'.} \\
\texttt{{-} (7,6) je '◌'.} \\
\texttt{{-} (7,7) je '◌'.} \\
\\ 
\texttt{Predlažem da koristimo lijevi stupac (stupac 1) i prvi redak (redak 1) jer su tamo linije minimalne.} \\
\texttt{Cilj:} \\
\texttt{{-} C na (1,1) {-}> Ali L je na (1,1). Moramo prvo premjestiti L.} \\
\texttt{{-} L na (1,2) {-}> (1,2) je prazno.} \\
\texttt{{-} P na (1,3) {-}> (1,3) je prazno.} \\
\\ 
\texttt{Ali (1,1) je zauzeto L{-}om. Dakle, prvo moramo premjestiti L s (1,1) na (1,2).} \\
\texttt{Zatim C s (5,5) na (1,1).} \\
\texttt{Zatim P s (7,2) na (1,3).} \\
\\ 
\texttt{Međutim, (1,2) je prazno, ali (1,1) je zauzeto L{-}om.} \\
\texttt{Dakle, prvi korak: Premjesti L s (1,1) na (1,2).} \\
\\ 
\texttt{Ali čekaj, (1,2) je prazno? Da, na početnoj ploči (1,2) je '◌'.} \\
\\ 
\texttt{Dakle, predlažem:} \\
\texttt{1. L ide na (1,2)} \\
\texttt{2. C ide na (1,1)} \\
\texttt{3. P ide na (1,3)} \\
\\ 
\texttt{Molim te, potvrdi da su (1,2), (1,1) i (1,3) prazne na tvojoj ploči. Ako jesu, možemo nastaviti.} \\
\\ 
\texttt{Ali (1,1) je zauzeto L{-}om na mojoj ploči. Dakle, (1,1) nije prazno dok ne premjestim L.} \\
\texttt{Dakle, redoslijed je:} \\
\texttt{1. Premjesti L s (1,1) na (1,2).} \\
\texttt{2. Premjesti C s (5,5) na (1,1).} \\
\texttt{3. Premjesti P s (7,2) na (1,3).} \\
\\ 
\texttt{Molim te, provjeri jesu li (1,2), (1,1) i (1,3) prazne na tvojoj ploči. Ako jesu, možemo pokušati.} \\
\\ 
\texttt{Ali (1,1) je zauzeto L{-}om na mojoj ploči, pa (1,1) nije prazno dok ne premjestim L.} \\
\texttt{Dakle, (1,1) će biti prazno tek nakon što premjestim L.} \\
\\ 
\texttt{Dakle, predlažem da prvo premjestim L na (1,2).} \\
\\ 
\texttt{POMAKNI: L, (1, 2)} \\
            }
        }
    }
     \\ \\

    \theutterance \stepcounter{utterance}  
    & & & \multicolumn{2}{p{0.3\linewidth}}{
        \cellcolor[rgb]{0.9,0.9,0.9}{
            \makecell[{{p{\linewidth}}}]{
                \texttt{\tiny{[GM$|$GM]}}
                \texttt{RECI: Već imamo 6 kazni, što je kritično. Nažalost, i (3,2) je zauzeto linijom. Čini se da su mnoge ćelije u sredini i desnom dijelu blokirane linijama. Moramo biti vrlo oprezni.} \\
\\ 
\texttt{Pogledajmo ponovno moju ploču:} \\
\texttt{{-} L je na (1,1)} \\
\texttt{{-} C je na (5,5)} \\
\texttt{{-} P je na (7,2)} \\
\\ 
\texttt{Linije se nalaze na:} \\
\texttt{{-} (2,2) do (2,4) (gornji dio)} \\
\texttt{{-} (4,4) do (4,5) (sredina)} \\
\texttt{{-} (5,4) do (5,5) (sredina)} \\
\texttt{{-} (6,4) do (6,5) (donji dio)} \\
\\ 
\texttt{Ćelije koje su sigurno prazne i blizu su:} \\
\texttt{{-} (1,2) je prazno? Ne, (1,2) je '◌' na početku, ali provjerimo. Na početnoj ploči, (1,2) je '◌'.} \\
\texttt{{-} (1,3) je '◌'.} \\
\texttt{{-} (1,4) je '◌'.} \\
\texttt{{-} (2,1) je '◌'.} \\
\texttt{{-} (3,1) je '◌'.} \\
\texttt{{-} (3,3) je '◌'.} \\
\texttt{{-} (4,1) je '◌'.} \\
\texttt{{-} (4,2) je '◌'.} \\
\texttt{{-} (4,3) je '◌'.} \\
\texttt{{-} (5,1) je '◌'.} \\
\texttt{{-} (5,2) je '◌'.} \\
\texttt{{-} (5,3) je '◌'.} \\
\texttt{{-} (6,1) je '◌'.} \\
\texttt{{-} (6,2) je '◌'.} \\
\texttt{{-} (6,3) je '◌'.} \\
\texttt{{-} (7,1) je '◌'.} \\
\texttt{{-} (7,3) je '◌'.} \\
\texttt{{-} (7,4) je '◌'.} \\
\texttt{{-} (7,5) je '◌'.} \\
\texttt{{-} (7,6) je '◌'.} \\
\texttt{{-} (7,7) je '◌'.} \\
\\ 
\texttt{Predlažem da koristimo lijevi stupac (stupac 1) i prvi redak (redak 1) jer su tamo linije minimalne.} \\
\texttt{Cilj:} \\
\texttt{{-} C na (1,1) {-}> Ali L je na (1,1). Moramo prvo premjestiti L.} \\
\texttt{{-} L na (1,2) {-}> (1,2) je prazno.} \\
\texttt{{-} P na (1,3) {-}> (1,3) je prazno.} \\
\\ 
\texttt{Ali (1,1) je zauzeto L{-}om. Dakle, prvo moramo premjestiti L s (1,1) na (1,2).} \\
\texttt{Zatim C s (5,5) na (1,1).} \\
\texttt{Zatim P s (7,2) na (1,3).} \\
\\ 
\texttt{Međutim, (1,2) je prazno, ali (1,1) je zauzeto L{-}om.} \\
\texttt{Dakle, prvi korak: Premjesti L s (1,1) na (1,2).} \\
\\ 
\texttt{Ali čekaj, (1,2) je prazno? Da, na početnoj ploči (1,2) je '◌'.} \\
\\ 
\texttt{Dakle, predlažem:} \\
\texttt{1. L ide na (1,2)} \\
\texttt{2. C ide na (1,1)} \\
\texttt{3. P ide na (1,3)} \\
\\ 
\texttt{Molim te, potvrdi da su (1,2), (1,1) i (1,3) prazne na tvojoj ploči. Ako jesu, možemo nastaviti.} \\
\\ 
\texttt{Ali (1,1) je zauzeto L{-}om na mojoj ploči. Dakle, (1,1) nije prazno dok ne premjestim L.} \\
\texttt{Dakle, redoslijed je:} \\
\texttt{1. Premjesti L s (1,1) na (1,2).} \\
\texttt{2. Premjesti C s (5,5) na (1,1).} \\
\texttt{3. Premjesti P s (7,2) na (1,3).} \\
\\ 
\texttt{Molim te, provjeri jesu li (1,2), (1,1) i (1,3) prazne na tvojoj ploči. Ako jesu, možemo pokušati.} \\
\\ 
\texttt{Ali (1,1) je zauzeto L{-}om na mojoj ploči, pa (1,1) nije prazno dok ne premjestim L.} \\
\texttt{Dakle, (1,1) će biti prazno tek nakon što premjestim L.} \\
\\ 
\texttt{Dakle, predlažem da prvo premjestim L na (1,2).} \\
\\ 
\texttt{POMAKNI: L, (1, 2)} \\
            }
        }
    }
    & & \\ \\

    \theutterance \stepcounter{utterance}  
    & & & \multicolumn{2}{p{0.3\linewidth}}{
        \cellcolor[rgb]{0.9,0.9,0.9}{
            \makecell[{{p{\linewidth}}}]{
                \texttt{\tiny{[GM$|$GM]}}
                \texttt{Invalid response format: RECI: Već imamo 6 kazni, što je kritično. Nažalost, i (3,2) je zauzeto linijom. Čini se da su mnoge ćelije u sredini i desnom dijelu blokirane linijama. Moramo biti vrlo oprezni.} \\
\\ 
\texttt{Pogledajmo ponovno moju ploču:} \\
\texttt{{-} L je na (1,1)} \\
\texttt{{-} C je na (5,5)} \\
\texttt{{-} P je na (7,2)} \\
\\ 
\texttt{Linije se nalaze na:} \\
\texttt{{-} (2,2) do (2,4) (gornji dio)} \\
\texttt{{-} (4,4) do (4,5) (sredina)} \\
\texttt{{-} (5,4) do (5,5) (sredina)} \\
\texttt{{-} (6,4) do (6,5) (donji dio)} \\
\\ 
\texttt{Ćelije koje su sigurno prazne i blizu su:} \\
\texttt{{-} (1,2) je prazno? Ne, (1,2) je '◌' na početku, ali provjerimo. Na početnoj ploči, (1,2) je '◌'.} \\
\texttt{{-} (1,3) je '◌'.} \\
\texttt{{-} (1,4) je '◌'.} \\
\texttt{{-} (2,1) je '◌'.} \\
\texttt{{-} (3,1) je '◌'.} \\
\texttt{{-} (3,3) je '◌'.} \\
\texttt{{-} (4,1) je '◌'.} \\
\texttt{{-} (4,2) je '◌'.} \\
\texttt{{-} (4,3) je '◌'.} \\
\texttt{{-} (5,1) je '◌'.} \\
\texttt{{-} (5,2) je '◌'.} \\
\texttt{{-} (5,3) je '◌'.} \\
\texttt{{-} (6,1) je '◌'.} \\
\texttt{{-} (6,2) je '◌'.} \\
\texttt{{-} (6,3) je '◌'.} \\
\texttt{{-} (7,1) je '◌'.} \\
\texttt{{-} (7,3) je '◌'.} \\
\texttt{{-} (7,4) je '◌'.} \\
\texttt{{-} (7,5) je '◌'.} \\
\texttt{{-} (7,6) je '◌'.} \\
\texttt{{-} (7,7) je '◌'.} \\
\\ 
\texttt{Predlažem da koristimo lijevi stupac (stupac 1) i prvi redak (redak 1) jer su tamo linije minimalne.} \\
\texttt{Cilj:} \\
\texttt{{-} C na (1,1) {-}> Ali L je na (1,1). Moramo prvo premjestiti L.} \\
\texttt{{-} L na (1,2) {-}> (1,2) je prazno.} \\
\texttt{{-} P na (1,3) {-}> (1,3) je prazno.} \\
\\ 
\texttt{Ali (1,1) je zauzeto L{-}om. Dakle, prvo moramo premjestiti L s (1,1) na (1,2).} \\
\texttt{Zatim C s (5,5) na (1,1).} \\
\texttt{Zatim P s (7,2) na (1,3).} \\
\\ 
\texttt{Međutim, (1,2) je prazno, ali (1,1) je zauzeto L{-}om.} \\
\texttt{Dakle, prvi korak: Premjesti L s (1,1) na (1,2).} \\
\\ 
\texttt{Ali čekaj, (1,2) je prazno? Da, na početnoj ploči (1,2) je '◌'.} \\
\\ 
\texttt{Dakle, predlažem:} \\
\texttt{1. L ide na (1,2)} \\
\texttt{2. C ide na (1,1)} \\
\texttt{3. P ide na (1,3)} \\
\\ 
\texttt{Molim te, potvrdi da su (1,2), (1,1) i (1,3) prazne na tvojoj ploči. Ako jesu, možemo nastaviti.} \\
\\ 
\texttt{Ali (1,1) je zauzeto L{-}om na mojoj ploči. Dakle, (1,1) nije prazno dok ne premjestim L.} \\
\texttt{Dakle, redoslijed je:} \\
\texttt{1. Premjesti L s (1,1) na (1,2).} \\
\texttt{2. Premjesti C s (5,5) na (1,1).} \\
\texttt{3. Premjesti P s (7,2) na (1,3).} \\
\\ 
\texttt{Molim te, provjeri jesu li (1,2), (1,1) i (1,3) prazne na tvojoj ploči. Ako jesu, možemo pokušati.} \\
\\ 
\texttt{Ali (1,1) je zauzeto L{-}om na mojoj ploči, pa (1,1) nije prazno dok ne premjestim L.} \\
\texttt{Dakle, (1,1) će biti prazno tek nakon što premjestim L.} \\
\\ 
\texttt{Dakle, predlažem da prvo premjestim L na (1,2).} \\
\\ 
\texttt{POMAKNI: L, (1, 2)} \\
            }
        }
    }
    & & \\ \\

    \theutterance \stepcounter{utterance}  
    & & \multicolumn{4}{p{0.6\linewidth}}{
        \cellcolor[rgb]{0.9,0.9,0.9}{
            \makecell[{{p{\linewidth}}}]{
                \texttt{\tiny{[GM$\rangle$P2]}}
                \texttt{Kazna: Tvoja poruka sadrži više od jedne naredbe!} \\
\\ 
\texttt{Provjerite da vaš odgovor sadrži samo `RECI: <PORUKA>` ili `POMAKNI: <OBJEKT>, (<X>, <Y>)`, i ništa drugo!} \\
\\ 
\texttt{Zajedno ste prikupili 7 od 8 kazni.} \\
\texttt{Molim pokušajte ponovno!} \\
            }
        }
    }
    & \\ \\

    \theutterance \stepcounter{utterance}  
    & & & \multicolumn{4}{p{0.6\linewidth}}{
        \cellcolor[rgb]{0.9,0.9,0.9}{
            \makecell[{{p{\linewidth}}}]{
                \texttt{\tiny{[GM$\langle$P2]}}
                \texttt{RECI: Imamo samo još jednu kaznu. Moramo biti precizni. Na mojoj ploči, L je na (1,1), C na (5,5), P na (7,2). Predlažem da L premjestimo na (1,2), C na (1,1), P na (1,3). Provjeri jesu li (1,2), (1,1) i (1,3) prazne na tvojoj ploči. Ako jesu, potvrdi i ja ću pokušati premjestiti L na (1,2).} \\
            }
        }
    }
     \\ \\

    \theutterance \stepcounter{utterance}  
    & & & \multicolumn{2}{p{0.3\linewidth}}{
        \cellcolor[rgb]{0.9,0.9,0.9}{
            \makecell[{{p{\linewidth}}}]{
                \texttt{\tiny{[GM$|$GM]}}
                \texttt{RECI: Imamo samo još jednu kaznu. Moramo biti precizni. Na mojoj ploči, L je na (1,1), C na (5,5), P na (7,2). Predlažem da L premjestimo na (1,2), C na (1,1), P na (1,3). Provjeri jesu li (1,2), (1,1) i (1,3) prazne na tvojoj ploči. Ako jesu, potvrdi i ja ću pokušati premjestiti L na (1,2).} \\
            }
        }
    }
    & & \\ \\

    \theutterance \stepcounter{utterance}  
    & & \multicolumn{4}{p{0.6\linewidth}}{
        \cellcolor[rgb]{0.9,0.9,0.9}{
            \makecell[{{p{\linewidth}}}]{
                \texttt{\tiny{[P1$\langle$GM]}}
                \texttt{Vaša poruka je proslijeđena drugom igraču.} \\
\\ 
\\ 
\texttt{Trenutačno igrate rundu 4 od 12.} \\
\texttt{Zajedno ste prikupili 7 od 8 kazni.} \\
\texttt{Drugi igrač je poslao ovu poruku:} \\
\\ 
\\ 
\texttt{"Imamo samo još jednu kaznu. Moramo biti precizni. Na mojoj ploči, L je na (1,1), C na (5,5), P na (7,2). Predlažem da L premjestimo na (1,2), C na (1,1), P na (1,3). Provjeri jesu li (1,2), (1,1) i (1,3) prazne na tvojoj ploči. Ako jesu, potvrdi i ja ću pokušati premjestiti L na (1,2)."} \\
\\ 
\\ 
\texttt{Koja je vaša sljedeća naredba?} \\
            }
        }
    }
    & \\ \\

    \theutterance \stepcounter{utterance}  
    & \multicolumn{4}{p{0.6\linewidth}}{
        \cellcolor[rgb]{0.9,0.9,0.9}{
            \makecell[{{p{\linewidth}}}]{
                \texttt{\tiny{[P1$\rangle$GM]}}
                \texttt{RECI: Nažalost, (1,2) i (1,3) nisu prazne na mojoj ploči jer tamo postoje linije. Tvoja prethodna strategija s kolonom 3 je bila bolja jer su (3,1), (3,2) i (3,3) prazne. Molim te, potvrdi da se vraćamo na dogovor: C na (3,1), L na (3,2), P na (3,3). Ja ću sada pokušati premjestiti L na (3,2).} \\
            }
        }
    }
    & & \\ \\

    \theutterance \stepcounter{utterance}  
    & & & \multicolumn{2}{p{0.3\linewidth}}{
        \cellcolor[rgb]{0.9,0.9,0.9}{
            \makecell[{{p{\linewidth}}}]{
                \texttt{\tiny{[GM$|$GM]}}
                \texttt{RECI: Nažalost, (1,2) i (1,3) nisu prazne na mojoj ploči jer tamo postoje linije. Tvoja prethodna strategija s kolonom 3 je bila bolja jer su (3,1), (3,2) i (3,3) prazne. Molim te, potvrdi da se vraćamo na dogovor: C na (3,1), L na (3,2), P na (3,3). Ja ću sada pokušati premjestiti L na (3,2).} \\
            }
        }
    }
    & & \\ \\

    \theutterance \stepcounter{utterance}  
    & & \multicolumn{4}{p{0.6\linewidth}}{
        \cellcolor[rgb]{0.9,0.9,0.9}{
            \makecell[{{p{\linewidth}}}]{
                \texttt{\tiny{[GM$\rangle$P2]}}
                \texttt{Vaša poruka je proslijeđena drugom igraču.} \\
\\ 
\\ 
\texttt{Trenutačno igrate rundu 5 od 12.} \\
\texttt{Zajedno ste prikupili 7 od 8 kazni.} \\
\texttt{Drugi igrač je poslao ovu poruku:} \\
\\ 
\\ 
\texttt{"Nažalost, (1,2) i (1,3) nisu prazne na mojoj ploči jer tamo postoje linije. Tvoja prethodna strategija s kolonom 3 je bila bolja jer su (3,1), (3,2) i (3,3) prazne. Molim te, potvrdi da se vraćamo na dogovor: C na (3,1), L na (3,2), P na (3,3). Ja ću sada pokušati premjestiti L na (3,2)."} \\
\\ 
\\ 
\texttt{Koja je vaša sljedeća naredba?} \\
            }
        }
    }
    & \\ \\

    \theutterance \stepcounter{utterance}  
    & & & \multicolumn{4}{p{0.6\linewidth}}{
        \cellcolor[rgb]{0.9,0.9,0.9}{
            \makecell[{{p{\linewidth}}}]{
                \texttt{\tiny{[GM$\langle$P2]}}
                \texttt{POMAKNI: L, (3, 2)} \\
            }
        }
    }
     \\ \\

    \theutterance \stepcounter{utterance}  
    & & & \multicolumn{2}{p{0.3\linewidth}}{
        \cellcolor[rgb]{0.9,0.9,0.9}{
            \makecell[{{p{\linewidth}}}]{
                \texttt{\tiny{[GM$|$GM]}}
                \texttt{POMAKNI: L, (3, 2)} \\
            }
        }
    }
    & & \\ \\

    \theutterance \stepcounter{utterance}  
    & & & \multicolumn{2}{p{0.3\linewidth}}{
        \cellcolor[rgb]{0.9,0.9,0.9}{
            \makecell[{{p{\linewidth}}}]{
                \texttt{\tiny{[GM$|$GM]}}
                \texttt{Kazna: (3,2) nije prazno, već sadrži '─'.} \\
\\ 
\texttt{Zajedno ste prikupili 8 od 8 kazni.} \\
\texttt{Molim pokušajte ponovno!} \\
            }
        }
    }
    & & \\ \\

    \theutterance \stepcounter{utterance}  
    & & \multicolumn{4}{p{0.6\linewidth}}{
        \cellcolor[rgb]{0.9,0.9,0.9}{
            \makecell[{{p{\linewidth}}}]{
                \texttt{\tiny{[GM$\rangle$P2]}}
                \texttt{Kazna: (3,2) nije prazno, već sadrži '─'.} \\
\\ 
\texttt{Zajedno ste prikupili 8 od 8 kazni.} \\
\texttt{Molim pokušajte ponovno!} \\
            }
        }
    }
    & \\ \\

    \theutterance \stepcounter{utterance}  
    & & & \multicolumn{4}{p{0.6\linewidth}}{
        \cellcolor[rgb]{0.9,0.9,0.9}{
            \makecell[{{p{\linewidth}}}]{
                \texttt{\tiny{[GM$\langle$P2]}}
                \texttt{RECI: gotovo?} \\
            }
        }
    }
     \\ \\

    \theutterance \stepcounter{utterance}  
    & & & \multicolumn{2}{p{0.3\linewidth}}{
        \cellcolor[rgb]{0.9,0.9,0.9}{
            \makecell[{{p{\linewidth}}}]{
                \texttt{\tiny{[GM$|$GM]}}
                \texttt{RECI: gotovo?} \\
            }
        }
    }
    & & \\ \\

    \theutterance \stepcounter{utterance}  
    & & \multicolumn{4}{p{0.6\linewidth}}{
        \cellcolor[rgb]{0.9,0.9,0.9}{
            \makecell[{{p{\linewidth}}}]{
                \texttt{\tiny{[P1$\langle$GM]}}
                \texttt{Vaša poruka je proslijeđena drugom igraču.} \\
\\ 
\\ 
\texttt{Trenutačno igrate rundu 5 od 12.} \\
\texttt{Zajedno ste prikupili 8 od 8 kazni.} \\
\texttt{Drugi igrač je poslao ovu poruku:} \\
\\ 
\\ 
\texttt{"gotovo?"} \\
\\ 
\\ 
\texttt{Koja je vaša sljedeća naredba?} \\
            }
        }
    }
    & \\ \\

    \theutterance \stepcounter{utterance}  
    & \multicolumn{4}{p{0.6\linewidth}}{
        \cellcolor[rgb]{0.9,0.9,0.9}{
            \makecell[{{p{\linewidth}}}]{
                \texttt{\tiny{[P1$\rangle$GM]}}
                \texttt{RECI: gotovo!} \\
            }
        }
    }
    & & \\ \\

    \theutterance \stepcounter{utterance}  
    & & & \multicolumn{2}{p{0.3\linewidth}}{
        \cellcolor[rgb]{0.9,0.9,0.9}{
            \makecell[{{p{\linewidth}}}]{
                \texttt{\tiny{[GM$|$GM]}}
                \texttt{RECI: gotovo!} \\
            }
        }
    }
    & & \\ \\

    \theutterance \stepcounter{utterance}  
    & & \multicolumn{4}{p{0.6\linewidth}}{
        \cellcolor[rgb]{0.9,0.9,0.9}{
            \makecell[{{p{\linewidth}}}]{
                \texttt{\tiny{[GM$\rangle$P2]}}
                \texttt{Vaša poruka je proslijeđena drugom igraču.} \\
\\ 
\\ 
\texttt{Trenutačno igrate rundu 6 od 12.} \\
\texttt{Zajedno ste prikupili 8 od 8 kazni.} \\
\texttt{Drugi igrač je poslao ovu poruku:} \\
\\ 
\\ 
\texttt{"gotovo!"} \\
\\ 
\\ 
\texttt{Koja je vaša sljedeća naredba?} \\
            }
        }
    }
    & \\ \\

    \theutterance \stepcounter{utterance}  
    & & & \multicolumn{4}{p{0.6\linewidth}}{
        \cellcolor[rgb]{0.9,0.9,0.9}{
            \makecell[{{p{\linewidth}}}]{
                \texttt{\tiny{[GM$\langle$P2]}}
                \texttt{RECI: gotovo!} \\
            }
        }
    }
     \\ \\

    \theutterance \stepcounter{utterance}  
    & & & \multicolumn{2}{p{0.3\linewidth}}{
        \cellcolor[rgb]{0.9,0.9,0.9}{
            \makecell[{{p{\linewidth}}}]{
                \texttt{\tiny{[GM$|$GM]}}
                \texttt{RECI: gotovo!} \\
            }
        }
    }
    & & \\ \\

    \theutterance \stepcounter{utterance}  
    & & & \multicolumn{2}{p{0.3\linewidth}}{
        \cellcolor[rgb]{0.9,0.9,0.9}{
            \makecell[{{p{\linewidth}}}]{
                \texttt{\tiny{[GM$|$GM]}}
                \texttt{true} \\
            }
        }
    }
    & & \\ \\

    \theutterance \stepcounter{utterance}  
    & & & \multicolumn{2}{p{0.3\linewidth}}{
        \cellcolor[rgb]{0.9,0.9,0.9}{
            \makecell[{{p{\linewidth}}}]{
                \texttt{\tiny{[GM$|$GM]}}
                \texttt{Initial states:} \\
\texttt{```} \\
\texttt{Player 1:} \\
\texttt{ 1234567} \\
\texttt{╔═╤═╤═╤═╗} \\
\texttt{║◌│◌│◌│◌║ 1} \\
\texttt{╟─┼─┤◌│◌║ 2} \\
\texttt{║◌│◌│◌│◌║ 3} \\
\texttt{║L│◌├─┴─╢ 4} \\
\texttt{╟─┴─┤◌◌◌║ 5} \\
\texttt{║◌◌C└─┬─╢ 6} \\
\texttt{║◌◌◌◌P│◌║ 7} \\
\texttt{╚═════╧═╝} \\
\\ 
\\ 
\texttt{Player 2:} \\
\texttt{ 1234567} \\
\texttt{╔═╤═╤═╤═╗} \\
\texttt{║L│◌│◌│◌║ 1} \\
\texttt{╟─┼─┤◌│◌║ 2} \\
\texttt{║◌│◌│◌│◌║ 3} \\
\texttt{║◌│◌├─┴─╢ 4} \\
\texttt{╟─┴─┤C◌◌║ 5} \\
\texttt{║◌◌◌└─┬─╢ 6} \\
\texttt{║◌P◌◌◌│◌║ 7} \\
\texttt{╚═════╧═╝} \\
\\ 
\texttt{```} \\
            }
        }
    }
    & & \\ \\

    \theutterance \stepcounter{utterance}  
    & & & \multicolumn{2}{p{0.3\linewidth}}{
        \cellcolor[rgb]{0.9,0.9,0.9}{
            \makecell[{{p{\linewidth}}}]{
                \texttt{\tiny{[GM$|$GM]}}
                \texttt{End states:} \\
\texttt{```} \\
\texttt{Player 1:} \\
\texttt{ 1234567} \\
\texttt{╔═╤═╤═╤═╗} \\
\texttt{║◌│◌│◌│◌║ 1} \\
\texttt{╟─┼─┤◌│◌║ 2} \\
\texttt{║◌│◌│◌│◌║ 3} \\
\texttt{║L│◌├─┴─╢ 4} \\
\texttt{╟─┴─┤◌◌◌║ 5} \\
\texttt{║◌◌C└─┬─╢ 6} \\
\texttt{║◌◌◌◌P│◌║ 7} \\
\texttt{╚═════╧═╝} \\
\\ 
\\ 
\texttt{Player 2:} \\
\texttt{ 1234567} \\
\texttt{╔═╤═╤═╤═╗} \\
\texttt{║L│◌│◌│◌║ 1} \\
\texttt{╟─┼─┤◌│◌║ 2} \\
\texttt{║◌│◌│◌│◌║ 3} \\
\texttt{║◌│◌├─┴─╢ 4} \\
\texttt{╟─┴─┤C◌◌║ 5} \\
\texttt{║◌◌◌└─┬─╢ 6} \\
\texttt{║◌P◌◌◌│◌║ 7} \\
\texttt{╚═════╧═╝} \\
\\ 
\texttt{```} \\
            }
        }
    }
    & & \\ \\

    \theutterance \stepcounter{utterance}  
    & & & \multicolumn{2}{p{0.3\linewidth}}{
        \cellcolor[rgb]{0.9,0.9,0.9}{
            \makecell[{{p{\linewidth}}}]{
                \texttt{\tiny{[GM$|$GM]}}
                \texttt{* success: True} \\
\texttt{* lose: False} \\
\texttt{* aborted: False} \\
\texttt{{-}{-}{-}{-}{-}{-}{-}} \\
\texttt{* Shifts: 0.00} \\
\texttt{* Max Shifts: 4.00} \\
\texttt{* Min Shifts: 2.00} \\
\texttt{* End Distance Sum: 8.24} \\
\texttt{* Init Distance Sum: 8.24} \\
\texttt{* Expected Distance Sum: 12.57} \\
\texttt{* Penalties: 8.00} \\
\texttt{* Max Penalties: 8.00} \\
\texttt{* Rounds: 6.00} \\
\texttt{* Max Rounds: 12.00} \\
\texttt{* Object Count: 3.00} \\
            }
        }
    }
    & & \\ \\

\end{supertabular}
}

\end{document}
